\documentclass[a4paper, serif, 10pt]{moderncv}

%% moderncv
% color and style
\moderncvcolor{black}
\moderncvstyle{banking}

% renew moderncv command to adapt for CJK colon
\renewcommand*{\cvitem}[3][.25em]{%
  \ifstrempty{#2}{}{\hintstyle{#2}：}{#3}%
  \par\addvspace{#1}}

\renewcommand*{\cvitemwithcomment}[4][.25em]{%
  \savebox{\cvitemwithcommentmainbox}{\ifstrempty{#2}{}{\hintstyle{#2}：}#3}%
  \setlength{\cvitemwithcommentmainlength}{\widthof{\usebox{\cvitemwithcommentmainbox}}}%
  \setlength{\cvitemwithcommentcommentlength}{\maincolumnwidth-\separatorcolumnwidth-\cvitemwithcommentmainlength}%
  \begin{minipage}[t]{\cvitemwithcommentmainlength}\usebox{\cvitemwithcommentmainbox}\end{minipage}%
  \hfill% fill of \separatorcolumnwidth
  \begin{minipage}[t]{\cvitemwithcommentcommentlength}\raggedleft\small\itshape#4\end{minipage}%
  \par\addvspace{#1}}

%% page layout/margins
\usepackage[top=2.5cm, bottom=2.5cm, left=1.5cm, right=1.5cm]{geometry}

\nopagenumbers{}

%% Babel config for English language
\usepackage[english]{babel}

%% fontspec
\usepackage{fontspec}

\IfFontExistsTF{Linux Libertine}{
  \setmainfont[Ligatures={TeX, Common}, Numbers=Lining, ItalicFont=Linux Libertine]{Linux Libertine}
}{}
\IfFontExistsTF{Linux Libertine O}{
  \setmainfont[Ligatures={TeX, Common}, Numbers=Lining, ItalicFont=Linux Libertine O]{Linux Libertine O}
}{}

%% CTeX
% CJK support, used to show CJK characters in the resume
%
% - heading=false: disable ctex heading
% - punct=kaiming: use kaiming punctuations styles for CJK
% - scheme=plain: use plain scheme, do not override \normalsize font size
% - space=auto: space settings for CJK characters
%
% CJK fonts are resolved per locale language, e.g. Noto Serif/Sans CJK JP
% for Japanese resumes, with fallback to the next region if not installed
%
% ref:
% - http://ctan.mirrorcatalogs.com/language/chinese/ctex/ctex.pdf
\usepackage[UTF8, heading=false, punct=kaiming, scheme=plain, space=auto]{ctex}

\IfFontExistsTF{Noto Serif CJK HK}{
  \setCJKmainfont{Noto Serif CJK HK}
}{
\IfFontExistsTF{Noto Serif CJK TC}{
  \setCJKmainfont{Noto Serif CJK TC}
}{}}
\IfFontExistsTF{Noto Sans CJK HK}{
  \setCJKsansfont{Noto Sans CJK HK}
}{
\IfFontExistsTF{Noto Sans CJK TC}{
  \setCJKsansfont{Noto Sans CJK TC}
}{}}

%% line spacing
\usepackage{setspace}
\setstretch{1.125}

%% URL styling - use normal text instead of monospace
\urlstyle{same}

%% Patch \httplink and \httpslink to handle full URLs
% The original moderncv macros always prepend http:// or https:// to URLs.
% This causes issues when the URL already has a protocol (e.g., https://example.com)
% resulting in malformed URLs like https://https://example.com.
% This patch checks if the URL already contains "://" and uses it directly if so.
% Note: We use \str_set:Nx (x-expansion) to expand macros like \@homepage first.
\ExplSyntaxOn
\str_new:N \g_yamlresume_url_str

\RenewDocumentCommand{\httpslink}{O{}m}{
  \str_set:Nx \g_yamlresume_url_str {#2}
  \str_if_in:NnTF \g_yamlresume_url_str {://}
    { \url{#2} }
    { \url{https://#2} }
}

\RenewDocumentCommand{\httplink}{O{}m}{
  \str_set:Nx \g_yamlresume_url_str {#2}
  \str_if_in:NnTF \g_yamlresume_url_str {://}
    { \url{#2} }
    { \url{http://#2} }
}
\ExplSyntaxOff

\name{陳嘉欣}{}
\title{人力資源經理｜人才發展與組織效能}
\phone[mobile]{+852 9123 4567}
\email{ka.yan.chan@example.com}
\homepage{https://www.chan-kayan-hr.com}

\address{中國香港，香港島，香港，香港中環皇后大道中99號}{}{}

\extrainfo{{\small \faLinkedin}\ \href{https://www.linkedin.com/in/chan-kayan-hr}{@chan-kayan-hr} {} {} {} • {} {} {} 
{\small \faMedium}\ \href{https://medium.com/@chan-kayan-hr}{@chan-kayan-hr}}

\begin{document}

\maketitle

\section{簡介}

\cvline{}{具備超過十年香港及大中華區人力資源管理經驗，專注人才策略、組織發展與員工體驗。擅長在快速變化的零售及科技環境中，將業務目標轉化為可執行的人才方案。
%
\begin{itemize}
\item 曾領導 8 人人力資源團隊，支持逾 1,200 名員工，並將年度流失率降低 18\%。
\item 主導薪酬架構重整及績效管理制度優化，提升員工投入度分數 12 個百分點。
\item 熟悉香港僱傭條例、強積金、稅務及跨境用工合規，並具備勞資協商經驗。
\item 以數據驅動人力規劃，建立招聘漏斗、離職分析及人才盤點儀表板。
\end{itemize}}
\section{教育背景}

\cventry{2012年9月–2014年6月}
        {碩士，人力資源管理，成績：3.8}
        {香港中文大學}
        {\url{https://www.cuhk.edu.hk}}
        {}
        {\begin{itemize}
\item 深入研究人才策略與組織效能，並以零售業個案完成畢業論文
\item 透過小組項目建立薪酬設計、人才盤點及變革管理實務能力
\item 獲選參與學院與企業合辦的人力資源顧問項目
\end{itemize}
\textbf{課程}：策略人力資源管理、組織行為學、薪酬與績效管理、招聘與選才、培訓與發展、勞動法規與僱傭實務}

\cventry{2008年9月–2012年6月}
        {學士，心理學，成績：3.6}
        {香港城市大學}
        {\url{https://www.cityu.edu.hk}}
        {}
        {\begin{itemize}
\item 以心理學角度理解員工行為、團隊動力與組織文化
\item 完成關於工作動機與員工投入度的專題研究
\item 培養數據分析、訪談及問卷設計的基礎能力
\end{itemize}
\textbf{課程}：工業與組織心理學、研究方法與統計、消費者行為、社會心理學、溝通與人際關係}
\section{工作經歷}

\cventry{2021年4月至今}
        {人力資源經理}
        {宏基零售集團有限公司}
        {\url{https://www.hongkei-retail.com}}
        {}
        {\begin{itemize}
\item 帶領 8 人人力資源團隊，負責招聘、薪酬、員工關係、培訓及人力資源資訊系統
\item 與管理層制定三年人才策略，支持 12 間門店及逾 1,200 名員工的業務擴展
\item 重整薪酬架構與年度調薪機制，將薪酬預算準確度提升至 98\%
\item 建立月度人力數據儀表板，追蹤流失率、招聘周期及員工投入度
\item 處理複雜勞資個案及紀律調查，確保符合香港僱傭條例及公司政策
\end{itemize}
\textbf{關鍵字}：人才策略、組織發展、員工關係、績效管理、薪酬架構}

\cventry{2017年7月–2021年3月}
        {高級人力資源主任}
        {卓越科技顧問有限公司}
        {\url{https://www.excellence-tech.com.hk}}
        {}
        {\begin{itemize}
\item 負責香港及深圳兩地辦公室的全流程招聘，每年處理逾 150 個職位
\item 設計校園招聘及實習生計劃，將應屆畢業生留任率提升 20\%
\item 協助推行績效管理制度，培訓 30 位部門主管進行目標設定與回饋面談
\item 與財務部門合作優化薪酬預算及福利方案，控制年度人力成本增長
\item 建立員工手冊及入職流程，提升新員工首 90 天體驗
\end{itemize}
\textbf{關鍵字}：全流程招聘、校園招聘、績效管理、員工手冊、跨境用工}

\cventry{2014年7月–2017年6月}
        {人力資源主任}
        {香港人力資源服務有限公司}
        {\url{https://www.hkhrservices.com}}
        {}
        {\begin{itemize}
\item 為中小企客戶提供招聘、薪酬計算及僱傭合規支援
\item 處理每月薪酬、強積金、報稅及假期管理，確保準時及準確
\item 協助客戶編寫職位說明書及設計面試評分表
\item 回應員工查詢並支援勞資糾紛初步調解
\item 定期更新香港僱傭條例及市場薪酬資訊予客戶
\end{itemize}
\textbf{關鍵字}：薪酬計算、強積金、僱傭合規、職位說明書、客戶支援}
\section{語言}

\cvline{粵語}{母語或雙語水平 \hfill \textbf{關鍵字}：母語、香港粵語、商務溝通}
\cvline{普通話}{完全專業水平 \hfill \textbf{關鍵字}：普通話水平測試二級甲等、跨部門溝通}
\cvline{英語}{完全專業水平 \hfill \textbf{關鍵字}：雅思 7.5、商務書寫、簡報匯報}
\section{專業技能}

\cvline{人力資源策略}{專家 \hfill \textbf{關鍵字}：人才規劃、組織設計、繼任計劃、變革管理}
\cvline{招聘與僱主品牌}{高級 \hfill \textbf{關鍵字}：全流程招聘、校園招聘、僱主品牌、面試評估}
\cvline{薪酬與福利}{高級 \hfill \textbf{關鍵字}：薪酬架構、市場薪酬調研、福利設計、薪酬預算}
\cvline{員工關係與合規}{高級 \hfill \textbf{關鍵字}：勞資協商、香港僱傭條例、員工手冊、紀律調查}
\cvline{學習與發展}{高級 \hfill \textbf{關鍵字}：培訓需求分析、課程設計、領導力發展、學習成效評估}
\cvline{人力資源資訊系統}{中級 \hfill \textbf{關鍵字}：人力數據分析、儀表板設計、報表自動化、流程數碼化}
\section{榮譽}

\cventry{2023年11月}
        {香港人力資源管理學會}
        {年度卓越人力資源領袖獎}
        {}
        {}
        {表揚在人才策略、員工體驗及組織變革方面具傑出表現的人力資源領袖。}

\cventry{2022年9月}
        {亞洲人力資源大獎}
        {最佳僱主品牌項目金獎}
        {}
        {}
        {憑校園招聘品牌重塑及員工推薦計劃，提升僱主知名度及年輕人才申請量。}
\section{證書}

\cventry{2019年6月}
        {人力資源認證協會}
        {全球人力資源專業人士認證（GPHR）}
        {\url{https://www.hrci.org}}
        {}
        {}

\cventry{2020年11月}
        {香港和解中心}
        {認可調解員}
        {\url{https://www.mediationcentre.org.hk}}
        {}
        {}
\section{出版刊物}

\cventry{2023年8月}
        {以員工體驗驅動人才保留：香港零售業實證研究}
        {香港人力資源管理學會月刊}
        {\url{https://www.hkihrm.org}}
        {}
        {\begin{itemize}
\item 分析零售業前線員工流失的關鍵因素
\item 提出結合入職體驗、前線主管培訓及彈性排班的留才方案
\item 以三家香港零售企業的試點數據說明員工投入度與留任率的關係
\end{itemize}}
\section{項目}

\cventry{2022年2月–2023年5月}
        {整合入職、學習、績效與離職數據，建立員工體驗儀表板。}
        {員工體驗數碼轉型計劃}
        {}
        {}
        {\begin{itemize}
\item 與資訊科技及營運部門合作，將五個人力資源流程數碼化
\item 建立離職風險指標，協助主管提前介入關鍵人才保留
\item 將入職文件處理時間由 3 天縮短至 1 天
\item 為 50 位主管提供數據解讀及員工關懷培訓
\end{itemize}
\textbf{關鍵字}：員工體驗、流程數碼化、人力數據、留才}

\cventry{2019年9月–2020年2月}
        {重新設計校園招聘宣傳、面試流程及實習生體驗。}
        {校園招聘品牌重塑}
        {}
        {}
        {\begin{itemize}
\item 與市場部合作製作僱主品牌內容，提升社交媒體互動率 35\%
\item 設計結構化面試評分表，統一不同部門的選才標準
\item 舉辦線上職業講座，吸引逾 600 名應屆畢業生參加
\item 實習生轉正率由 45\% 提升至 62\%
\end{itemize}
\textbf{關鍵字}：僱主品牌、校園招聘、結構化面試、實習生計劃}
\section{興趣愛好}

\cvline{運動}{馬拉松、瑜伽、行山}
\cvline{文化}{閱讀、攝影、旅行}
\section{志願服務}

\cventry{2018年9月至今}
        {職涯規劃義務顧問}
        {香港青年協會}
        {\url{https://www.hkfyg.org.hk}}
        {}
        {\begin{itemize}
\item 為應屆畢業生提供履歷指導及模擬面試
\item 主持職涯探索工作坊，協助青年了解人力資源及零售行業
\item 協助設計求職技巧教材，供中心義工重複使用
\item 每年服務超過 80 名青年
\end{itemize}}
    
\end{document}